\documentclass[../../main.tex]{subfiles}
\begin{document}
\section{Sard 定理}\label{section:ch06-Sard定理}

下面是支撑我们所有关于嵌入、逼近与横截性结果的定理。

\begin{theorem}[Sard's Theorem]\label{theorem:ch06-sard}
设 $M$ 与 $N$ 是带边或不带边光滑流形, $F:M\to N$ 是光滑映射。则 $F$ 的临界值之集在 $N$ 中零测。
\end{theorem}

\begin{proof}
令 $m=\dim M$, $n=\dim N$。我们对 $m$ 归纳。当 $m=0$ 时结论显然:若 $n=0$, $F$ 无临界点;若 $n>0$, $F$ 的整个像可数,故零测。

现设 $m\geqslant 1$,并假设定理对定义域维数小于 $m$ 的映射成立。用可数多个光滑坐标卡覆盖 $M$ 与 $N$ 后,可归结为 $F$ 是从 $U\subseteq\mathbb{R}^m$ 或 $\mathbb{H}^m$ 的开子集到 $\mathbb{R}^n$ 的光滑映射这一情形。把定义域 $U$ 中的坐标记为 $\bigl(x^1,\dots,x^m\bigr)$,陪域中的坐标记为 $\bigl(y^1,\dots,y^n\bigr)$。

令 $C\subseteq U$ 表示 $F$ 的临界点集。我们定义一列递减子集
\begin{align*}
C\supseteq C_1\supseteq C_2\supseteq\cdots
\end{align*}
如下:
\begin{align}\label{align:ch06-Ck}
C_k=\{x\in C:\text{对 }1\leqslant i\leqslant k,\ F\text{ 的所有 }i\text{ 阶偏导数在 }x\text{ 处为零}\}.
\end{align}
由连续性, $C$ 及所有 $C_k$ 在 $U$ 中闭。我们将分三步证明 $F(C)$ 零测。

\textbf{第一步}: $F(C\setminus C_1)$ 零测。由于 $C_1$ 在 $U$ 中闭,我们不妨把 $U$ 换成 $U\setminus C_1$,并假设 $C_1=\varnothing$。令 $a\in C$ 为一点。我们的假设意味着 $F$ 的某个一阶偏导数在 $a$ 处非零。

通过重排定义域与陪域中的坐标,可设 $\partial F^1/\partial x^1(a)\neq 0$。这意味着我们可在 $a$ 在 $U$ 中的某邻域 $V_a$ 内定义新的光滑坐标
\begin{align*}
(u,v)=\bigl(u,v^2,\dots,v^m\bigr)
\end{align*}
为
\begin{align*}
u=F^1,\quad v^2=x^2,\quad \dots,\quad v^m=x^m,
\end{align*}
因为该坐标变换的 Jacobi 矩阵在 $a$ 处非奇异。必要时缩小 $V_a$,可设 $\overline{V_a}$ 是 $U$ 的紧子集,且坐标光滑延拓到 $\overline{V_a}$ 上。在这些坐标下, $F$ 具有坐标表示
\begin{align}\label{align:ch06-F-coords}
F\bigl(u,v^2,\dots,v^m\bigr)=\bigl(u,F^2(u,v),\dots,F^n(u,v)\bigr),
\end{align}
其 Jacobi 矩阵为
\begin{align*}
\mathrm{D}F(u,v)=
\begin{pmatrix}
1 & 0\\
\dfrac{\partial F^i}{\partial v^j}
\end{pmatrix}.
\end{align*}
因此, $C\cap\overline{V_a}$ 恰好由使 $(n-1)\times(m-1)$ 矩阵 $\bigl(\partial F^i/\partial v^j\bigr)$ 的秩小于 $n-1$ 的点组成。

我们要证明集合 $F(C\cap\overline{V_a})$ 在 $\mathbb{R}^n$ 中零测。由于该集紧,由\cref{lemma:ch06-compact-hyperplane} 只需证明它与每个超平面 $y^1=c$ 的交具有 $(n-1)$ 维零测度。

对 $c\in\mathbb{R}$,令 $B_c=\{v:(c,v)\in\overline{V_a}\}\subseteq\mathbb{R}^{m-1}$,并定义 $F_c:B_c\to\mathbb{R}^{n-1}$ 为
\begin{align*}
F_c(v)=\bigl(F^2(c,v),\dots,F^n(c,v)\bigr).
\end{align*}
由于 $F(c,v)=(c,F_c(v))$, $F|_{\overline{V_a}}$ 落在超平面 $y^1=c$ 上的临界值恰为形如 $(c,w)$ 的点,其中 $w$ 是 $F_c$ 的临界值。由归纳假设,每个 $F_c$ 的临界值集具有 $(n-1)$ 维零测度。于是由\cref{lemma:ch06-compact-hyperplane}, $F(C\cap\overline{V_a})$ 零测。

由于 $U$ 被可数多个形如 $\overline{V_a}$ 的集合覆盖, $F(C\cap U)$ 是零测集的可数并,故零测。这就完成了第一步。

\textbf{第二步}: 对每个 $k$, $F(C_k\setminus C_{k+1})$ 零测。同样,由于 $C_{k+1}$ 在 $U$ 中闭,我们可丢弃它并假设在 $C_k$ 的每一点都有 $F$ 的某个 $(k+1)$ 阶偏导数非零。

任取 $a\in C_k$,并令 $\gamma:U\to\mathbb{R}$ 表示 $F$ 的某个 $k$ 阶偏导数,它在 $a$ 处至少有一个一阶偏导数非零。于是 $a$ 是光滑映射 $\gamma$ 的正则点,故存在 $a$ 的邻域 $V_a$,它完全由 $\gamma$ 的正则点组成。令 $Y$ 为 $\gamma$ 在 $V_a$ 中的零点集,由正则水平集定理它是光滑超曲面。由 $C_k$ 的定义, $F$ 的所有 $k$ 阶导数(包括 $\gamma$)在 $C_k$ 上为零,故 $C_k\cap V_a\subseteq Y$。在任一点 $p\in C_k\cap V_a$ 处, $\mathrm{d}F_p$ 不满射,当然 $\mathrm{d}(F|_Y)_p=(\mathrm{d}F_p)|_{T_pY}$ 也不满射。因此, $F(C_k\cap V_a)$ 含于 $F|_Y:Y\to\mathbb{R}^n$ 的临界值集中,而由归纳假设后者零测。由于 $U$ 可被可数多个形如 $V_a$ 的邻域覆盖, $F(C_k\setminus C_{k+1})$ 含于形如 $F(C_k\cap V_a)$ 的集合的可数并中,故零测。

我们尚未完成,因为 $C$ 中可能有这样的点, $F$ 的所有偏导数在其上都为零,这意味着它们既不在 $C\setminus C_1$ 中,也不在任何 $C_k\setminus C_{k+1}$ 中。最后一步处理这种可能性。

\textbf{第三步}: 对 $k>m/n-1$, $F(C_k)$ 零测。对每个 $a\in U$,存在闭立方体 $E$ 使 $a\in E\subseteq U$。由于 $U$ 可被可数多个这样的立方体覆盖,只需证明当 $E$ 是含于 $U$ 中的闭立方体时, $F(C_k\cap E)$ 零测。令 $E$ 为这样的立方体,并令 $A$ 为 $F$ 在 $E$ 中所有 $(k+1)$ 阶导数绝对值的上界常数。令 $R$ 表示 $E$ 的边长, $K$ 为待选的大整数。我们可把 $E$ 细分为 $K^m$ 个边长为 $R/K$ 的立方体,记为 $\{E_1,\dots,E_{K^m}\}$。

若 $E_i$ 是其中一个立方体且存在点 $a_i\in C_k\cap E_i$,则 Taylor 定理的\cref{corollary:appC-taylor} 表明:对所有 $x\in E_i$,有
\begin{align}\label{align:ch06-taylor-bound}
|F(x)-F(a_i)|\leqslant A'|x-a_i|^{k+1},
\end{align}
其中 $A'$ 是只依赖于 $A,k,m$ 的常数。于是 $F(E_i)$ 含于半径为 $A'(R/K)^{k+1}$ 的球中。这蕴涵 $F(C_k\cap E)$ 含于 $K^m$ 个球之并中,其 $n$ 维体积之和不超过
\begin{align}\label{align:ch06-volume-bound}
K^m\cdot(A')^n(R/K)^{n(k+1)}=A''K^{m-nk-n},
\end{align}
其中 $A''=(A')^nR^{n(k+1)}$。由于 $k>m/n-1$,通过取 $K$ 充分大可使上式任意小,故 $F(C_k\cap E)$ 零测。
\end{proof}

\begin{corollary}\label{corollary:ch06-image-zero-dim}
设 $M$ 与 $N$ 是带边或不带边光滑流形, $F:M\to N$ 是光滑映射。若 $\dim M<\dim N$,则 $F(M)$ 在 $N$ 中零测。
\end{corollary}

\begin{proof}
在此情形, $M$ 的每一点都是 $F$ 的临界点。
\end{proof}

习题 6-1 给出了上述推论的一个简单证明,它不依赖 Sard 定理的全部强度。

重要的是要意识到:若 $F$ 仅假设连续,则\cref{corollary:ch06-image-zero-dim} 不成立。例如,存在连续映射 $F:[0,1]\to\mathbb{R}^2$,其像是整个单位正方形 $[0,1]\times[0,1]$。(这种映射称为空间填充曲线。例子见 [Rud76, 第 168 页]。)

\begin{corollary}\label{corollary:ch06-immersed-submanifold-zero}
设 $M$ 是带边或不带边光滑流形, $S\subseteq M$ 是带边或不带边浸入子流形。若 $\dim S<\dim M$,则 $S$ 在 $M$ 中零测。
\end{corollary}

\begin{proof}
对包含映射 $S\hookrightarrow M$ 应用\cref{corollary:ch06-image-zero-dim}。
\end{proof}

\end{document}